% Working revision after close reading of the supplied ASC/ESWA papers.
% Text only: insert inside the existing abstract environment.
Repeated clustering determines how sensing and forwarding work are
distributed in energy-limited underwater wireless sensor networks.
In each collection round, selected cluster heads organize sensor reports
into routes toward a sink.
The task is to choose these heads as residual energies change, while
accounting for the consumption induced by the resulting routing tree.
Such decisions affect both network expenditure and the fractions of
remaining energy consumed at individual nodes.
However, ranking candidates by total energy alone can give similar scores
to trees that impose substantially different fractional drains on weak nodes.
This paper investigates a density-weighted fractional-depletion criterion
for repeated cluster-head selection in three-dimensional deployments.
The criterion combines modeled total consumption with a powered penalty
on node depletion, using calibrated reference scales to compare candidates
within a round.
We propose DACO, a density-aware ant colony optimization method that uses
remaining grid support during candidate construction and the proposed
criterion when ranking completed trees.
Shared tree construction and node-pheromone feedback provide the framework
for evaluating the effects of the density and depletion mechanisms.
\textbf{[Final experimental sentence required: state the supported findings,
the relevant deployment conditions, and the observed computational trade-off.]}
